% Vietnamese translation for OI Public Library
% Source-File: IOI中国国家候选队论文/2006/俞鑫/俞鑫.doc
\documentclass[11pt]{article}
\usepackage{oipl-vn}
\usepackage{tikz}

\title{Bàn cờ trong bàn cờ\\\large Bàn về tư tưởng phân hoạch bàn cờ}
\author{Yu Xin, Trường Trung học trực thuộc Đại học Phục Đán}
\date{2006}

\begin{document}
\maketitle

\origtitle{Bàn cờ trong bàn cờ: bàn về tư tưởng phân hoạch bàn cờ}
\sourcepath{IOI Candidate Team Papers/2006/Yu Xin/Yu Xin.doc}

\begin{abstract}
Các bài thi tin học hiện nay thường mượn một mô hình toán học nào đó làm môi
trường đặt đề, khiến bài toán vừa thú vị vừa có nền tảng toán học sâu. Bàn cờ
là một trong những mô hình toán học quan trọng như vậy. Bài viết tập trung
phân tích một tư tưởng quan trọng trên bàn cờ: tư tưởng phân hoạch bàn cờ. Qua
hai ví dụ điển hình, bài viết nêu các quy luật nên tuân theo khi phân hoạch
bàn cờ, để người đọc có một nhận thức nhất định về những cách phân hoạch bàn
cờ vốn rất đa dạng và phức tạp.
\end{abstract}

\noindent\textbf{Từ khóa:} mô hình toán học, bàn cờ, tư tưởng thuật toán.

\tableofcontents

\section{Dẫn nhập}

Tin học là một ngành tổng hợp, đồng thời cũng là một ngành đầy thú vị. Bàn cờ,
với tư cách một mô hình toán học quan trọng, thường được người ra đề ưa dùng
nhờ tính thú vị và các tính chất toán học phức tạp của nó. Vì vậy, nghiên cứu
sâu các tư tưởng thuật toán ẩn trong bàn cờ là việc rất cần thiết đối với
người yêu thích tin học. Ở đây, tôi sẽ tập trung trình bày một tư tưởng quan
trọng trong các bài toán bàn cờ: tư tưởng phân hoạch bàn cờ.

Với một bàn cờ \(m\times n\), số bàn cờ con hình chữ nhật đã rất lớn, chẳng
hạn theo cách đếm bằng biên ta có
\[
  \binom{m+1}{2}\binom{n+1}{2}
\]
bàn con. Số cách phân hoạch nó lại càng nhiều hơn. Nếu phân hoạch bàn cờ một
cách khéo léo, ta có thể giải được nhiều loại bài toán trên bàn cờ.

\section{Ví dụ 1: Phủ bàn cờ}

Đây là một bài toán kinh điển.

\subsection{Mô tả bài toán}

Trong một bàn cờ gồm \(2^k\times 2^k\) ô vuông, nếu đúng một ô khác với các ô
còn lại, ta gọi ô đó là \term{ô đặc biệt}, và gọi bàn cờ đó là \term{bàn cờ
đặc biệt}. Rõ ràng vị trí của ô đặc biệt trên bàn cờ có \(4^k\) khả năng. Vì
thế, với mọi \(k\ge 0\), có \(4^k\) bàn cờ đặc biệt khác nhau.

Trong bài toán phủ bàn cờ, ta cần dùng bốn dạng quân domino chữ L khác nhau để
phủ mọi ô không đặc biệt của một bàn cờ đặc biệt cho trước, và hai quân chữ L
bất kỳ không được phủ chồng lên nhau. Dễ thấy trong mọi cách phủ một bàn
\(2^k\times 2^k\), số quân chữ L cần dùng đúng bằng
\[
  \frac{4^k-1}{3}.
\]
Hãy tìm một cách phủ.

\paragraph{Dữ liệu vào.}
Dòng đầu tiên là \(k\), tức kích thước bàn cờ. Dòng thứ hai gồm \(x,y\),
\(1\le x,y\le 2^k\), lần lượt là hàng và cột của ô đặc biệt.

\paragraph{Dữ liệu ra.}
Gồm \(2^k\) dòng, mỗi dòng có \(2^k\) số, biểu thị số hiệu quân chữ L phủ ô
tương ứng. Ô đặc biệt được ghi là \(0\).

\paragraph{Ví dụ.}
\begin{verbatim}
Dữ liệu vào:
2
1 2

Dữ liệu ra:
1 0 2 2
1 1 3 2
4 3 3 5
4 4 5 5
\end{verbatim}

\subsection{Phân tích thuật toán}

Vì kích thước bàn cờ là \(2^k\times 2^k\), ta tự nhiên nghĩ tới việc chia nó
thành bốn bàn cờ con kích thước \(2^{k-1}\times 2^{k-1}\). Tuy nhiên, chỉ một
trong bốn bàn cờ con chứa ô đặc biệt; ba bàn cờ còn lại không phải bàn cờ đặc
biệt. Vì vậy vấn đề vẫn chưa được giải quyết.

Chỉ cần suy nghĩ thêm một chút, ta thấy rằng nếu đặt một quân chữ L ở chính
giữa bàn cờ, che ba ô trung tâm thuộc ba bàn cờ con không chứa ô đặc biệt, thì
cả bốn bàn cờ con đều trở thành bàn cờ đặc biệt. Lúc này bài toán biến thành
bốn bài toán con cùng dạng, nên có thể giải bằng đệ quy đơn giản.

\begin{figure}[ht]
\centering
\begin{tikzpicture}[scale=0.72, every node/.style={font=\small}]
\draw[step=1] (0,0) grid (4,4);
\draw[very thick] (2,0) -- (2,4);
\draw[very thick] (0,2) -- (4,2);
\fill[gray!20] (1,3) rectangle (2,4);
\node at (1.5,3.5) {\(\star\)};
\fill[blue!20] (2,2) rectangle (3,3);
\fill[blue!20] (1,1) rectangle (2,2);
\fill[blue!20] (2,1) rectangle (3,2);
\node at (2.5,2.5) {\(t\)};
\node at (1.5,1.5) {\(t\)};
\node at (2.5,1.5) {\(t\)};
\node[align=center] at (2,-0.65)
  {Ba ô trung tâm được xem là ô đặc biệt\\cho ba bàn con còn lại.};
\end{tikzpicture}
\caption{Phân hoạch bàn \(4\times4\): đặt một quân chữ L ở tâm để biến bốn
bàn con thành bốn bài toán cùng dạng.}
\end{figure}

Ta dùng các ký hiệu sau:
\begin{description}[leftmargin=5.5em,style=nextline]
  \item[\(\texttt{num}\)] Mảng hai chiều; \(\texttt{num}[x,y]\) là số hiệu
  quân chữ L phủ ô \((x,y)\). Khi nói gán giá trị cho một ô, ta gán giá trị
  cho phần tử tương ứng của \(\texttt{num}\).
  \item[\(x_1,y_1\)] Hàng và cột của ô góc trên bên trái của bàn cờ hiện tại.
  \item[\(x_2,y_2\)] Hàng và cột của ô góc dưới bên phải của bàn cờ hiện tại.
  \item[\(x_3,y_3\)] Hàng và cột của ô đặc biệt trong bàn cờ hiện tại.
  \item[\(c_k\)] Kích thước hiện tại, tức bàn cờ có dạng
  \(2^{c_k}\times 2^{c_k}\).
  \item[\(\texttt{cnum}\)] Số hiệu quân chữ L hiện tại.
\end{description}

Ban đầu đặt \(\texttt{num}[x,y]=0\). Nếu \(c_k=0\), kích thước bàn cờ là
\(1\times 1\); ô duy nhất đã là ô đặc biệt được gán giá trị, nên không cần làm
gì. Nếu \(c_k>0\), đặt
\[
  x_m=\frac{x_1+x_2+1}{2},\qquad
  y_m=\frac{y_1+y_2+1}{2}.
\]
So sánh \(x_3\) với \(x_m\), và \(y_3\) với \(y_m\), ta biết ô đặc biệt nằm
trong bàn cờ con nào. Với ba bàn cờ con còn lại, gán ba ô sát tâm bàn cờ bằng
\(\texttt{cnum}\), đồng thời xem ba ô đó là ô đặc biệt của ba bàn cờ con ấy.
Sau đó tăng \(\texttt{cnum}\) lên \(1\), rồi đệ quy xử lý bốn bàn cờ con.

\subsection{Phân tích độ phức tạp}

Độ phức tạp thời gian là \(O(4^k)\), và độ phức tạp không gian cũng là
\(O(4^k)\). Vì số quân chữ L cần để phủ một bàn
\(2^k\times 2^k\) là \((4^k-1)/3\), thuật toán này là tối ưu theo nghĩa tiệm
cận.

\paragraph{Chương trình tham khảo.}
\texttt{GAMERS.PAS}

\subsection{Tiểu kết}

Khi chia bàn cờ thành các bàn cờ con, nên tuân theo hai điểm sau:
\begin{enumerate}
  \item Các bàn cờ được tách ra nên giống bàn cờ ban đầu nhiều nhất có thể.
  \item Sau khi chia bàn cờ ban đầu, cố gắng không để lại phần dư.
\end{enumerate}

Nhưng nếu sau khi chia chắc chắn phải còn phần dư thì sao? Ví dụ tiếp theo trả
lời câu hỏi đó.

\section{Ví dụ 2: Bài toán cờ Khổng Minh}

Bài toán lấy từ URAL 1051.

\subsection{Mô tả bài toán}

Trên các nút của một bàn cờ vô hạn có một số quân cờ. Các quân cờ này tạo
thành một hình chữ nhật \(m\times n\), trong đó \(m\) là chiều cao,
\(n\) là chiều rộng, và \(1\le m,n\le 1000\). Ta có thể dùng một quân cờ nhảy
qua một quân cờ kề nó; quân bị nhảy qua sẽ bị lấy đi. Hãy tìm số quân cờ ít
nhất có thể còn lại.

\paragraph{Dữ liệu vào.}
\(m,n\).

\paragraph{Dữ liệu ra.}
Số quân cờ ít nhất còn lại.

\paragraph{Ví dụ.}
\begin{verbatim}
Dữ liệu vào:
3 4

Dữ liệu ra:
2
\end{verbatim}

\subsection{Phân tích thuật toán}

Trường hợp \(m=1\) hoặc \(n=1\) khá đặc biệt, nên trước hết ta xử lý trường hợp
\(m,n\ge 2\). Để tiện trình bày, ta gọi bàn cờ gồm đúng các vị trí đang có quân
là \term{bàn cờ thật}. Từ ví dụ, có thể thấy rằng với ba quân liên tiếp ở các
ô \(4,5,6\) trong một khối \(3\times 3\), nếu các ô \(1,2,3\) trống còn các ô
\(7,8,9\) đều có quân, thì có thể dùng một chuỗi thao tác để lấy đi ba quân ở
\(4,5,6\). Ta gọi ba quân \(4,5,6\) là \term{module 1}.

Tuy nhiên, sau vài thử nghiệm ta thấy chỉ dùng module 1 để phân hoạch bàn cờ
thật \(m\times n\) không hiệu quả. Nguyên nhân là module 1 mỗi lần xử lý đồng
thời ba hàng liên tiếp; khi \(m\) không chia hết cho \(3\), sau khi phân hoạch
sẽ luôn còn phần dư. Cần chú ý rằng một phần được dùng làm module 1 chỉ thực
sự trở thành module 1 sau khi các quân ở bên trái nó đã bị loại bỏ.

\begin{figure}[ht]
\centering
\begin{tikzpicture}[scale=0.72, every node/.style={font=\small}]
\draw[step=1] (0,0) grid (3,3);
\foreach \x/\y/\n in {
  0.5/2.5/1,1.5/2.5/2,2.5/2.5/3,
  0.5/1.5/4,1.5/1.5/5,2.5/1.5/6,
  0.5/0.5/7,1.5/0.5/8,2.5/0.5/9} {
  \node at (\x,\y) {\(\n\)};
}
\foreach \x in {0.5,1.5,2.5} {
  \draw[fill=gray!20] (\x,1.5) circle (0.18);
  \draw[fill=black] (\x,0.5) circle (0.18);
}
\node[above] at (1.5,3) {trống};
\node[right] at (3.15,1.5) {module 1};
\node[right] at (3.15,0.5) {có quân};
\end{tikzpicture}
\caption{Module 1: trong khối \(3\times3\), ba quân ở hàng giữa có thể bị
loại bỏ khi phía trên đã trống và phía dưới còn quân hỗ trợ.}
\end{figure}

Vì thế, ta cần xử lý phần dư. Khi \(m\) không chia hết cho \(3\), phần dư luôn
là \(1\) hoặc \(2\) hàng. Một hàng quân rất khó loại bỏ. Nếu chỉ còn \(1\)
hàng, do \(m\ge 2\) và \(m\bmod 3=1\), ta có \(m\ge 4\). Khi đó ta xem bốn
hàng trên cùng là phần dư; phần còn lại có \(m-4\) hàng, chia hết cho \(3\),
nên có thể phân hoạch bằng module 1. Bốn hàng quân lại có thể chia thành hai
phần, mỗi phần gồm hai hàng. Vì vậy, xử lý trường hợp hai hàng quân trở thành
điểm then chốt.

Qua thử nghiệm, ta thấy rằng với sáu quân ở các ô \(2,3,5,6,8,9\), nếu các ô
\(1,4,7\) trống và các ô \(11,12\) có quân, thì có thể dùng một chuỗi thao tác
để loại bỏ sáu quân đó. Ta gọi sáu quân này là \term{module 2}.

\begin{figure}[ht]
\centering
\begin{tikzpicture}[scale=0.72, every node/.style={font=\small}]
\draw[step=1] (0,0) grid (3,4);
\foreach \x/\y/\n in {
  0.5/3.5/1,1.5/3.5/2,2.5/3.5/3,
  0.5/2.5/4,1.5/2.5/5,2.5/2.5/6,
  0.5/1.5/7,1.5/1.5/8,2.5/1.5/9,
  0.5/0.5/10,1.5/0.5/11,2.5/0.5/12} {
  \node at (\x,\y) {\(\n\)};
}
\foreach \x/\y in {1.5/3.5,2.5/3.5,1.5/2.5,2.5/2.5,1.5/1.5,2.5/1.5} {
  \draw[fill=gray!20] (\x,\y) circle (0.18);
}
\foreach \x/\y in {1.5/0.5,2.5/0.5} {
  \draw[fill=black] (\x,\y) circle (0.18);
}
\node[left] at (-0.1,2.5) {trống};
\node[right] at (3.15,2.5) {module 2};
\node[right] at (3.15,0.5) {hỗ trợ};
\end{tikzpicture}
\caption{Module 2 xử lý phần dư hai hàng: sáu quân ở các ô
\(2,3,5,6,8,9\) được loại bỏ nhờ các ô trống bên trái và hai quân hỗ trợ
phía dưới.}
\end{figure}

Với hai công cụ module 1 và module 2, việc phân hoạch bàn cờ thật \(m\times n\)
trở nên thuận lợi. Với mọi bàn cờ thật \(m\times n\), khi \(n\ge 5\), xét bàn
cờ \(m\times 3\) tạo bởi ba cột ngoài cùng bên trái; ta loại bỏ nó như sau.

\begin{enumerate}
  \item Nếu \(m\) chia hết cho \(3\), trước hết chia ba cột trái thành
  \(m/3\) bàn cờ con \(3\times 3\). Sau đó chia mỗi bàn \(3\times 3\) thành ba
  bàn \(3\times 1\). Mỗi lần thao tác trên bàn \(3\times 1\) phía trên cùng và
  bên trái nhất; vì phía trái của nó không có quân, nó chắc chắn là một module
  1 và có thể bị loại bỏ.
  \item Nếu \(m\bmod 3=1\), do \(m\ge 2\) nên \(m\ge 4\). Bàn \(2\times 3\)
  trên cùng của bàn \(m\times 3\) là một module 2, loại bỏ nó. Bàn
  \(2\times 3\) trên cùng của bàn \((m-2)\times 3\) cũng là một module 2, lại
  loại bỏ nó. Phần còn lại là bàn \((m-4)\times 3\); vì \(m-4\) chia hết cho
  \(3\), ta xử lý như trường hợp 1.
  \item Nếu \(m\bmod 3=2\), bàn \(2\times 3\) trên cùng của bàn \(m\times 3\)
  là một module 2, loại bỏ nó. Phần còn lại là bàn \((m-2)\times 3\); vì
  \(m-2\) chia hết cho \(3\), ta xử lý như trường hợp 1.
\end{enumerate}

Giả sử \(m\) đồng dư với \(p\) modulo \(3\), \(n\) đồng dư với \(q\) modulo
\(3\), trong đó ta chọn \(2\le p,q\le 4\). Cụ thể, phần dư \(0,1,2\) lần lượt
được biểu diễn bởi \(3,4,2\). Với mọi bàn cờ thật \(m\times n\), khi
\(n\ge 5\), liên tục loại bỏ ba cột bên trái theo cách trên, kích thước bàn cờ
thật biến đổi như sau:
\[
  m\times n \to m\times(n-3) \to m\times(n-6) \to \cdots \to m\times q.
\]
Sau đó xoay bàn cờ theo chiều kim đồng hồ \(90^\circ\), kích thước trở thành
\(q\times m\), rồi tiếp tục:
\[
  q\times m \to q\times(m-3) \to q\times(m-6) \to \cdots \to q\times p.
\]
Xoay thêm \(90^\circ\), ta còn bàn cờ \(p\times q\). Chỉ cần xử lý hữu hạn các
trường hợp \(p\times q\), ta thu được cận trên cho số quân ít nhất còn lại của
bàn cờ thật \(m\times n\).

\begin{center}
\begin{tabular}{c|ccc}
  \(\min s\) trên bàn \(p\times q\) & \(q=2\) & \(q=3\) & \(q=4\) \\
  \hline
  \(p=2\) & \(1\) & \(2\) & \(1\) \\
  \(p=3\) & \(2\) & \(2\) & \(2\) \\
  \(p=4\) & \(1\) & \(2\) & \(1\)
\end{tabular}
\end{center}

Trong các trường hợp có \(p=2,q=4\), \(p=3,q=3\), và \(p=4,q=4\), các module
1 hoặc module 2 còn lại được dùng để loại bỏ thêm quân bằng các biến thể xoay
hoặc lặp của hai mẫu ở trên. Trường hợp \(p=3,q=4\) chính là ví dụ đã nêu,
còn các trường hợp đối xứng nhận được bằng cách xoay bàn cờ \(90^\circ\).

Vì khi \(m\bmod 3=0,1,2\) thì \(p=3,4,2\) tương ứng, và khi
\(n\bmod 3=0,1,2\) thì \(q=3,4,2\) tương ứng, ta có thể viết cận trên theo
\(m,n\) như sau:
\[
\begin{array}{c|ccc}
  & n\bmod 3=0 & n\bmod 3=1 & n\bmod 3=2 \\
  \hline
  m\bmod 3=0 & 2 & 2 & 2 \\
  m\bmod 3=1 & 2 & 1 & 1 \\
  m\bmod 3=2 & 2 & 1 & 1
\end{array}
\]

Để tiện trình bày, gọi \(s\) là số quân ít nhất còn lại sau khi thao tác trên
bàn cờ thật \(m\times n\). Quan sát trên cho thấy: nếu trong \(m,n\) có một số
chia hết cho \(3\), thì cận trên của \(s\) là \(2\); nếu không, cận trên là
\(1\). Bây giờ ta chỉ cần chứng minh cận trên đó đúng là đáp án.

Thứ nhất, với mọi số nguyên dương \(m,n\), rõ ràng \(s\ge 0\). Nếu \(s=0\), vì
mỗi thao tác đúng bằng giảm một quân, nên ngay trước thao tác cuối cùng phải
còn đúng một quân. Nhưng với một quân duy nhất thì không thể thao tác, mâu
thuẫn. Do đó \(s\ge 1\).

Thứ hai, khi trong \(m,n\) có một số chia hết cho \(3\), không mất tính tổng
quát giả sử \(m\) chia hết cho \(3\); nếu không, xoay bàn cờ \(90^\circ\) để
đưa về trường hợp đó. Ta sẽ chứng minh \(s\ge 2\).

Tô màu xen kẽ bàn cờ vô hạn bằng ba màu sao cho trong mọi ba ô liên tiếp theo
hàng ngang hoặc hàng dọc đều có đúng một ô đỏ, một ô xanh lam, và một ô trắng.
Mỗi cột của bàn cờ thật \(m\times n\) là một bàn \(m\times 1\), và mỗi bàn
\(m\times 1\) lại chia được thành \(m/3\) bàn \(3\times 1\). Vì vậy bàn cờ
thật \(m\times n\) chia thành nhiều bàn \(3\times 1\), mỗi bàn có đúng một ô
đỏ, một ô trắng, và một ô xanh lam. Số ô đỏ, trắng, xanh lam trong bàn cờ thật
là như nhau. Gọi số quân trên ba màu đó lần lượt là \(a,b,c\); ban đầu
\(a=b=c\).

Trong mỗi thao tác, có thể xem như trên ba ô liên tiếp, hai ô mất quân còn ô
thứ ba thêm một quân. Vì vậy trong \(a,b,c\), có hai số giảm \(1\) và một số
tăng \(1\); tính chẵn lẻ của cả ba số đều đổi. Do đó \(a,b,c\) luôn cùng chẵn
lẻ. Nếu \(s=1\), sau khi thao tác xong ta có \(s=a+b+c\), tức trong
\(a,b,c\) có hai số bằng \(0\) và một số bằng \(1\); tính chẵn lẻ không giống
nhau, mâu thuẫn. Kết hợp với \(s\ge 1\), suy ra \(s\ge 2\).

Như vậy, với \(m,n\ge 2\), nếu một trong \(m,n\) chia hết cho \(3\) thì
\(s=2\), ngược lại \(s=1\).

\subsection{Trường hợp một hàng hoặc một cột}

Với \(m=1\), ta cũng trước hết dùng một ví dụ tổng quát để cho cận trên. Khi
\(m=1\), \(n\) quân cờ nằm trên một hàng; đánh số chúng từ trái sang phải là
\(1,2,3,\ldots,n\).

Đặt \(k=\lfloor n/2\rfloor\). Vì \(2k\le n\), ta thực hiện \(k\) thao tác.
Thao tác thứ \(i\), với \(1\le i\le k\), là cho quân số \(2i\) nhảy qua quân
số \(2i-1\). Trước thao tác thứ \(i\), các quân ở bên trái quân số \(2i-1\)
hoặc đã bị loại bỏ, hoặc đã di chuyển sang trái, còn các quân khác chưa di
chuyển. Vì vậy ô ngay bên trái quân số \(2i-1\) trống, nên thao tác hợp lệ.
Sau \(k\) thao tác, trên bàn còn \(n-\lfloor n/2\rfloor\) quân. Đây là cận
trên cho số quân còn lại ít nhất.

Bây giờ chỉ cần chứng minh \(n-\lfloor n/2\rfloor\) cũng là cận dưới. Trước
hết, nếu sau một thao tác nào đó không còn quân nào chưa từng tham gia thao
tác, thì về sau cũng sẽ không xuất hiện lại quân chưa từng tham gia thao tác.
Nếu có một thao tác như vậy, ta gọi nó là \term{thao tác then chốt}. Giai đoạn
có quân chưa từng tham gia thao tác kéo dài từ lúc bắt đầu đến thao tác then
chốt, hoặc đến khi không thể thao tác nữa.

Để chứng minh bài toán ban đầu, ta cần hai bổ đề. Xét cấu hình một hàng và gọi
các quân chưa từng tham gia thao tác là \term{quân chưa thao tác}. Ba điều kiện
sau sẽ được dùng:
\begin{enumerate}[label=\arabic*.,leftmargin=2.5em]
  \item Các quân chưa thao tác luôn xếp thành một đoạn liên tiếp.
  \item Ở bên trái đoạn quân chưa thao tác, hai quân đã thao tác bất kỳ không
  kề nhau; đồng thời hai ô ngay bên trái quân chưa thao tác ngoài cùng bên trái
  đều trống.
  \item Ở bên phải đoạn quân chưa thao tác, hai quân đã thao tác bất kỳ không
  kề nhau; đồng thời hai ô ngay bên phải quân chưa thao tác ngoài cùng bên
  phải đều trống.
\end{enumerate}

\paragraph{Bổ đề 1.}
Nếu một cấu hình đồng thời thỏa ba điều kiện trên, thì bước tiếp theo nếu có
thao tác sẽ phải dùng hai quân chưa thao tác.

\paragraph{Chứng minh.}
Xét một quân đã thao tác \(u\) ở bên trái đoạn quân chưa thao tác. Hai phía của
\(u\) hoặc là quân đã thao tác ở bên trái đoạn đó, hoặc là quân chưa thao tác
ngoài cùng bên trái. Theo điều kiện 2, giữa ô của \(u\) và các ô đó luôn có ít
nhất một ô trống, nên không thể thao tác với \(u\). Lập luận tương tự cho mọi
quân đã thao tác ở bên phải đoạn quân chưa thao tác. Do đó bước tiếp theo chỉ
có thể dùng hai quân chưa thao tác.

\paragraph{Bổ đề 2.}
Khi trên bàn còn quân chưa thao tác, phân bố quân luôn đồng thời thỏa ba điều
kiện trên.

\paragraph{Chứng minh.}
Ban đầu ba điều kiện đều đúng. Giả sử sau một thao tác nào đó, trên bàn vẫn còn
quân chưa thao tác nhưng cấu hình không còn đồng thời thỏa ba điều kiện; lấy
thao tác đầu tiên gây ra tình huống này. Trước thao tác đó, ba điều kiện đều
đúng; theo Bổ đề 1, hai quân được dùng trong thao tác ấy đều là quân chưa thao
tác.

Vì các quân chưa thao tác xếp thành một hàng liên tiếp, thao tác chỉ có hai
dạng: quân thứ hai từ trái nhảy qua quân thứ nhất từ trái, hoặc quân thứ hai
từ phải nhảy qua quân thứ nhất từ phải. Hai dạng đối xứng, nên chỉ xét dạng
thứ nhất. Gọi quân thứ nhất từ trái là \(a\), quân thứ hai là \(b\). Phía bên
phải không đổi nên điều kiện 3 vẫn đúng. Sau thao tác, trong đoạn quân chưa
thao tác ban đầu chỉ có hai quân ngoài cùng bên trái không còn là quân chưa
thao tác; các quân còn lại vẫn liên tiếp, nên điều kiện 1 đúng. Quân \(a\) bị
lấy đi, còn \(b\) nằm ở ô trống ngay bên trái vị trí cũ của \(a\). Vì trước đó
hai ô bên trái \(a\) đều trống, nên \(b\) không kề với các quân đã thao tác
khác ở bên trái đoạn quân chưa thao tác. Đồng thời hai ô bên trái quân chưa
thao tác ngoài cùng bên trái mới chính là hai ô cũ của \(a\) và \(b\), đều
trống. Vậy điều kiện 2 cũng đúng, mâu thuẫn với cách chọn thao tác đầu tiên.
Bổ đề được chứng minh.

Từ hai bổ đề suy ra hệ quả: khi còn quân chưa thao tác, bước tiếp theo nếu có
thao tác phải dùng hai quân chưa thao tác.

Ta chứng minh cận dưới bằng phản chứng. Giả sử
\(n-\lfloor n/2\rfloor\) không phải cận dưới, tức có thể dùng một chuỗi thao
tác để số quân còn lại nhỏ hơn \(n-\lfloor n/2\rfloor\). Vì mỗi thao tác giảm
đúng một quân, cần ít nhất \(\lfloor n/2\rfloor+1\) thao tác.

Mỗi thao tác nhiều nhất làm mất trạng thái ``chưa thao tác'' của hai quân. Do
đó trước bất kỳ thao tác nào trong \(\lfloor n/2\rfloor\) thao tác đầu tiên,
số quân chưa thao tác bị mất nhiều nhất là
\[
  2(\lfloor n/2\rfloor-1)<n.
\]
Vì thế vẫn còn ít nhất một quân chưa thao tác; theo hệ quả trên, thao tác phải
dùng hai quân chưa thao tác.

Nếu \(n\) lẻ, \(\lfloor n/2\rfloor=(n-1)/2\). Vậy \((n-1)/2\) thao tác đầu
đều dùng hai quân chưa thao tác. Sau các thao tác đó, đúng một quân chưa thao
tác còn lại; nhưng theo hệ quả, bước tiếp theo phải dùng hai quân chưa thao
tác, nên không thể thực hiện. Số thao tác nhỏ hơn
\(\lfloor n/2\rfloor+1\), mâu thuẫn.

Nếu \(n\) chẵn, \(\lfloor n/2\rfloor=n/2\). Khi đó \(n/2\) thao tác đầu đều
dùng hai quân chưa thao tác. Sau \(n/2-1\) thao tác đầu, đúng hai quân chưa
thao tác còn lại; theo hệ quả, thao tác tiếp theo phải dùng hai quân đó. Theo
Bổ đề 2, hai quân này kề nhau. Gọi quân bên trái là \(a'\), quân bên phải là
\(b'\). Hai khả năng \(a'\) nhảy qua \(b'\) và \(b'\) nhảy qua \(a'\) đối
xứng; xét \(a'\) nhảy qua \(b'\). Theo Bổ đề 2, các quân ở bên trái hai quân
này đôi một không kề nhau, và các quân ở bên phải cũng vậy. Do đó thao tác thứ
\(n/2+1\) nếu có bắt buộc phải dùng \(a'\). Sau thao tác thứ \(n/2\), quân
\(a'\) nằm ở ô trống bên phải \(b'\). Nhưng theo Bổ đề 2, hai ô bên phải
\(b'\) đều trống; vì vậy ô bên phải \(a'\) trống. Ô bên trái \(a'\) là vị trí
cũ của \(b'\), cũng trống sau khi \(b'\) bị lấy đi. Vậy \(a'\) không thể thao
tác, mâu thuẫn. Suy ra \(n-\lfloor n/2\rfloor\) là cận dưới.

Do đó khi \(m=1\), số quân ít nhất còn lại trên bàn cờ thật \(m\times n\) là
\(n-\lfloor n/2\rfloor\). Khi \(n=1\), xoay bàn \(m\times 1\) thành bàn
\(1\times m\), suy ra số quân ít nhất còn lại là
\(m-\lfloor m/2\rfloor\).

Tóm lại, số quân ít nhất còn lại là
\[
s=
\begin{cases}
  n-\lfloor n/2\rfloor, & m=1,\\
  m-\lfloor m/2\rfloor, & n=1,\\
  2, & m,n\ge 2\ \text{và}\ (3\mid m\ \text{hoặc}\ 3\mid n),\\
  1, & m,n\ge 2\ \text{và}\ 3\nmid m,\ 3\nmid n.
\end{cases}
\]

\paragraph{Chương trình tham khảo.}
\texttt{JUMP.PAS}

\subsection{Tiểu kết}

Trong bài toán này, ý tưởng chia bàn cờ thành các module 1 gặp trở ngại chủ
yếu vì sau khi chia luôn còn phần dư. Để xử lý phần dư, ta tìm ra module 2. Về
bản chất, module 2 chỉ là một số bước có quy luật dành riêng cho phần dư. Vì
vậy, khi việc phân hoạch bàn cờ gặp khó khăn và buộc phải để lại phần dư, có
thể dùng hai cách sau:
\begin{enumerate}
  \item Xử lý riêng phần dư, bằng một số bước đặc biệt ít ỏi, như xử lý bàn cờ
  nhỏ cuối cùng, hoặc bằng một số bước có quy luật, như module 2.
  \item Nếu phần dư thật sự khó xử lý, ta phải từ bỏ phương án phân hoạch ban
  đầu và tìm một phương án khác.
\end{enumerate}

\section{Tổng kết}

Là một tư tưởng quan trọng để xử lý bàn cờ, phương pháp phân hoạch bàn cờ thực
ra còn phong phú hơn rất nhiều. Về bản chất, tư tưởng phân hoạch bàn cờ là một
tư tưởng đệ quy, nhưng không hoàn toàn giống đệ quy thuần túy: nó phức tạp hơn
và khó nắm bắt hơn. Chỉ khi thường xuyên chú ý và suy nghĩ, ta mới có thể nắm
được nhiều điều tinh diệu trong đó, để khi thi đấu có thể ``quan sát rộng mà
chọn lọc gọn, tích lũy dày mà phát huy mạnh''.

\section*{Tài liệu tham khảo}

\begin{itemize}
  \item Wang Xiaodong, \textit{Thiết kế và phân tích thuật toán máy tính},
  tái bản lần 2, Nhà xuất bản Công nghiệp Điện tử.
  \item Ural State University Problem Set Archive.
\end{itemize}

\end{document}
